\section{Matrix operators}

\subsection{Transpose}
\index{matrices!transpose}

$\mat{A}\T$ denotes the transpose of $\mat{A}$, which flips the matrix across
its diagonal such that the rows become columns and vice versa.

\subsection{Inverse}
\index{matrices!inverse}

$\mat{A}^{-1}$ denotes the inverse of $\mat{A}$, which is defined as the matrix
such that $\mat{A}^{-1}\mat{A} = \mat{A}\mat{A}^{-1} = \mat{I}$.

\subsection{Inverse transpose}
\index{matrices!inverse transpose}

$\mat{A}^{-\mathsf{T}}$ denotes the inverse transpose of $\mat{A}$, which
performs an inverse then a transpose.

\subsection{Pseudoinverse}
\index{matrices!pseudoinverse}

$\mat{A}^+$ denotes the Moore-Penrose pseudoinverse of $\mat{A}$, which is a
generalization of the inverse matrix to nonsquare matrices. The pseudoinverse is
an approximate inverse in the least-squares sense, so it can solve least-squares
problems of the form $\mat{A}\mat{x} = \mat{b}$ via
$\mat{x} = \mat{A}^+\mat{b}$. See table \ref{tab:pseudoinverse} for
pseudoinverse definitions by system type.
\begin{booktable}
  \begin{tabular}{|cccc|}
    \hline
    \rowcolor{headingbg}
    \textbf{System type} & \textbf{Shape of A} & \textbf{Name of }$\mat{A}^+$ &
      \textbf{Value of }$\mat{A}^+$ \\
    \hline
    Well-determined & Square & Inverse &
      $\mat{A}^{-1}$ \\
    Overdetermined & Tall & Left pseudoinverse &
      $(\mat{A}\T\mat{A})^{-1}\mat{A}\T$ \\
    Underdetermined & Wide & Right pseudoinverse &
      $\mat{A}\T(\mat{A}\mat{A}\T)^{-1}$ \\
    \hline
  \end{tabular}
  \caption{Pseudoinverses by system type.}
  \label{tab:pseudoinverse}
\end{booktable}

\subsection{Diagonal}
\index{matrices!diagonal}

Let $\mat{x}$ be a vector and $\mat{A}$ and $\mat{B}$ be matrices.

$\diag(\mat{x})$ creates a diagonal matrix with the elements of $\mat{x}$ along
the diagonal.
\begin{equation*}
  \diag(\mat{x}) = \diag(x_1,\, \ldots,\, x_n) =
  \begin{bmatrix}
    x_1 & 0 & \cdots & 0 \\
    0 & x_2 & & \vdots \\
    \vdots & & \ddots & 0 \\
    0 & \cdots & 0 & x_n
  \end{bmatrix}
\end{equation*}

$\diag(\mat{A}, \mat{B})$ creates a block diagonal matrix.
\begin{equation*}
  \diag(\mat{A}, \mat{B}) =
  \begin{bmatrix}
    \mat{A} & \mat{0} \\
    \mat{0} & \mat{B}
  \end{bmatrix}
\end{equation*}

Operations on the $\diag()$ argument are applied element-wise.
\begin{equation*}
  \diag\left(\frac{1}{\mat{x}^2}\right) =
  \begin{bmatrix}
    \frac{1}{x_1^2} & 0 & \cdots & 0 \\
    0 & \frac{1}{x_2^2} & & \vdots \\
    \vdots & & \ddots & 0 \\
    0 & \cdots & 0 & \frac{1}{x_n^2}
  \end{bmatrix}
\end{equation*}

\subsection{Trace}
\index{matrices!trace}

$\tr(\mat{A})$ denotes the trace of the square matrix $\mat{A}$, which is
defined as the sum of the elements on the main diagonal.
